\documentclass[11pt]{article}
\usepackage[margin=1in]{geometry}
\usepackage{booktabs,longtable,array,graphicx,enumitem}
\usepackage[hidelinks]{hyperref}
\setlength{\parindent}{0pt}
\setlength{\parskip}{6pt}
\renewcommand{\arraystretch}{1.15}
\newcommand{\plainbox}[1]{\par\vspace{2pt}\noindent\fbox{\parbox{\dimexpr\linewidth-2\fboxsep-2\fboxrule\relax}{#1}}\par\vspace{2pt}}
\begin{document}
{\Large\textbf{Tasks Module, Part 1: The Questions}}\par
{\large What we ask, why we ask it, and where each question comes from}\par\vspace{4pt}\hrule\vspace{8pt}

\section*{1. What this part of the survey is for}
The ropeway may take away some of the work people do today. It may also bring new jobs.
To know who could move into a new job, we first need to know what work people really do,
and what work they have done before.

So we ask about \textbf{tasks}. A task is one small piece of work, such as ``carry a load'' or ``give change''.
Two jobs that share many tasks are close to each other. A person can move more easily
into a job that uses tasks they already know.

\section*{2. The rules we followed}
\begin{itemize}[leftmargin=1.2em,itemsep=2pt]
\item We ask what people \textbf{do}. We do not ask how good they are.
\item \textbf{No rating questions.} There are no scales and no numbers to pick. Answers are yes or no, a count, or a choice from a short list.
\item Questions are read aloud by the enumerator, with picture cards. Nobody needs to read.
\item Every question rests on a published survey or on the national list of occupations. Where a question has no published source, we say so.
\end{itemize}

\section*{3. How a task question works}
The enumerator reads each task and the person gives one of three answers:
\plainbox{\textbf{1} = Yes, I do this regularly in my main Yatra work.\par
\textbf{2} = Not in my main work, but I have done it before (another job, off-season work, or at home, for more than a few days in total).\par
\textbf{3} = I have never done it.}
Answer 2 matters. It shows skills that a job title hides. A porter who once worked on a building site
already knows how to build. Two other answers are recorded but not counted: \textbf{97} = does not know, and \textbf{98} = refused.

\section*{4. The 27 tasks}

Each task is read aloud with the same three-answer grid. The ``group'' column shows the broad group
we will use later in the analysis, and whether the work is done the same way each time (repeated) or
changes from case to case (varied). See section 7. The last column shows where the task comes from.

\textbf{Key to the short source names.} \emph{GS 2010} = Gathmann and Sch\"onberg (2010), a survey of 19 job tasks in Germany.
\emph{Spitz-Oener 2006} = a German study of job tasks over time. \emph{Autor-Levy-Murnane 2003} and \emph{Autor-Handel 2013} = United States task studies.
\emph{STEP} = the World Bank skills survey for developing countries (we hold the Philippines file). \emph{NCO} = India's National Classification of Occupations 2015, with the code and title of the job.
\emph{Pal et al. 2026} = a survey of Indian coal workers facing job change. \emph{Sengupta et al. 2021} = a survey of informal workers in Bangalore.
\emph{Apablaza et al. 2026} = a paper on measuring poor-quality work.
{\footnotesize
\begin{longtable}{p{0.7cm}p{4.2cm}p{2.7cm}p{6.9cm}}
\toprule
\textbf{No.} & \textbf{What we read aloud} & \textbf{Group} & \textbf{Where it comes from} \\
\midrule
\endfirsthead
\toprule
\textbf{No.} & \textbf{What we read aloud} & \textbf{Group} & \textbf{Where it comes from} \\
\midrule
\endhead
01 & Carry loads or luggage on foot & Hands-on / varied & \textbf{Adapted from} GS 2010 task 'pack, ship or transport'; NCO 9621.9900 luggage porters, 9621.0500 baggage porter. 'On foot on the trail' is NEW. \\
02 & Carry a passenger (palki, dandi, kandi) & Hands-on / varied & \textbf{New.} No NCO-2015 code for palki/dandi bearers (documented gap); nearest NCO 9331.0100 hand and pedal vehicle driver. Needs union check. \\
03 & Load, unload, stack or count goods & Hands-on / repeated & \textbf{Taken from} NCO 9333.0100 Loader and Unloader (shifting, stacking, counting loads). \\
04 & Pack or tie goods so they can be moved & Hands-on / repeated & \textbf{Taken from} GS 2010 task 'pack, ship or transport'; NCO 5223.0100 Shop Assistant (packs goods). \\
05 & Lead, load or control pack animals such as ponies or mules & Hands-on / varied & \textbf{Taken from} NCO 9332.0400 Pack Animal Driver (drives animal laden with goods, loads and unloads). \\
06 & Feed, groom or care for animals & Hands-on / varied & \textbf{Taken from} NCO 9332.0400 (feeds animals, cleans stables). \\
07 & Drive a motor vehicle (car, jeep, tempo, bus) & Hands-on / varied & \textbf{Taken from} NCO 8322.0501 Light Motor Vehicle Driver; Sengupta et al. 2021 skill 'driving'; Autor-Levy-Murnane 2003, Table I lists truck driving as a non-routine manual example. \\
08 & Operate an engine, machine or other equipment & Hands-on / repeated & \textbf{Taken from} GS 2010 'equip or operate machines'; Spitz-Oener 2006 routine manual; STEP module 6B 'operate heavy machines'. \\
09 & Repair or fix things by hand (vehicles, tools, saddles, buildings) & Hands-on / varied & \textbf{Taken from} GS 2010 'repair, renovate or reconstruct'; Spitz-Oener 2006 non-routine manual. \\
10 & Do electrical or wiring work & Hands-on / varied & \textbf{Taken from} NCO 7411.0100 Electrician, General (installs, maintains, repairs electrical equipment; NSQF level 4). \\
11 & Build or construct (masonry, carpentry, trail or road work) & Hands-on / varied & \textbf{Taken from} GS 2010 'manufacture, install or construct'; NCO 9313.9900 building construction labourers. \\
12 & Check equipment or loads for safety before use (saddle, rope, brakes) & Thinking and counting / repeated & \textbf{Adapted from} NCO 8343.1700 Aerial Ropeway Operator (checks cages clamped, rope in working order); Pal et al. 2026 task 'inspect machine and tools for potential faults'. \\
13 & Judge whether the weather or route is safe and decide whether to go & Thinking and counting / varied & \textbf{New.} No published counterpart. Idea anchored in the non-routine problem-solving task of Autor-Levy-Murnane 2003 and the 30-minute problem-solving item of Autor-Handel 2013. \\
14 & Sell goods or services to customers or pilgrims & Dealing with people / varied & \textbf{Taken from} GS 2010 'sell, buy or advertise'; Spitz-Oener 2006 'selling'; NCO 5223.0100 Shop Assistant, 4212.0100 Booking Clerk (issues tickets). \\
15 & Bargain or agree prices and fares & Dealing with people / varied & \textbf{Taken from} Spitz-Oener 2006 'negotiating'. \\
16 & Handle cash, give change or take digital payments & Thinking and counting / repeated & \textbf{Adapted from} NCO 4211.0100 Cashier (receives and counts cash); Sengupta et al. 2021 skill 'handling cash'. 'Digital payments' is NEW. \\
17 & Buy stock or arrange supplies from suppliers & Dealing with people / varied & \textbf{Taken from} GS 2010 'sell, buy'; Spitz-Oener 2006 'buying'; NCO 5151.0600 Hotel and Restaurant Keeper (purchases food stuff), 1120.2000 Retail proprietor (buys merchandise). \\
18 & Cook or prepare food or drink & Hands-on / varied & \textbf{Taken from} NCO 5120.0200 Cook, Institutional; 5212.9900 Street Food Vendors (Dhabawala). \\
19 & Serve guests (food, rooms, luggage) & Hands-on / varied & \textbf{Taken from} GS 2010 'serve or accommodate'; Spitz-Oener 2006 'serving or accommodating'; NCO 5131.0200 Steward, Hotel; 5151.0800 Room Bearer. \\
20 & Clean rooms, utensils or public areas & Hands-on / varied & \textbf{Taken from} GS 2010 'cleaning'; NCO 9112.9900 helpers and cleaners in hotels and establishments; Autor-Levy-Murnane 2003, Table I lists janitorial services as non-routine manual. \\
21 & Guide or explain routes, sites or rituals to visitors & Dealing with people / varied & \textbf{Taken from} NCO 5113.0200 Tourist Guide; Spitz-Oener 2006 'advising customers'. \\
22 & Give first aid or help sick or injured people & Hands-on / varied & \textbf{Adapted from} GS 2010 'nurse or treat others'. \\
23 & Keep order among crowds, queues or turns (baari), or watch over property & Hands-on / varied & \textbf{Adapted from} GS 2010 'secure'; NCO 5414.0501 Gateman / Unarmed Security Guard. Keeping a turn rotation (baari) is NEW. \\
24 & Coordinate with others by phone, radio or signals & Dealing with people / varied & \textbf{Taken from} NCO 8343.1700 Aerial Ropeway Operator (communicates by telephone or buzzer); GS 2010 'organise, coordinate'. \\
25 & Keep records or accounts (tickets, bookings, bills) & Thinking and counting / repeated & \textbf{Taken from} GS 2010 'calculate or do bookkeeping'; Spitz-Oener 2006 routine cognitive; NCO 4212.0100 Booking Clerk (maintains record of sale). \\
26 & Supervise, hire or organise other people's work & Dealing with people / varied & \textbf{Taken from} GS 2010 'employ, manage personnel, organise'; Spitz-Oener 2006 'employing or managing personnel'; STEP 'supervise others'. \\
27 & Teach or train others & Dealing with people / varied & \textbf{Taken from} GS 2010 'teach or train others'. \\
\bottomrule
\end{longtable}
}
\textbf{Main task.} After the list, one more question: ``Of the tasks you do regularly, which ONE takes most of your working time?'' The answer must be a task the person coded 1. Source: GS 2010: each worker also reports whether a task is his main activity; Autor-Handel 2013 asks about the share of the workday.

\section*{5. Questions about the job itself}
These eight questions describe how the work is done. None of them is a rating.
{\footnotesize
\begin{longtable}{p{0.7cm}p{4.5cm}p{3.0cm}p{6.4cm}}
\toprule
\textbf{No.} & \textbf{Question} & \textbf{Answers} & \textbf{Where it comes from} \\
\midrule
\endfirsthead
\toprule
\textbf{No.} & \textbf{Question} & \textbf{Answers} & \textbf{Where it comes from} \\
\midrule
\endhead
01 & In the busiest weeks, how many trips do you make in a day? (asked only if a carrying or animal task is coded 1) & whole number & \textbf{New.} Route-specific count. Replaces the STEP 1-10 physical-effort scale, which is a rating. Needs union check. \\
02 & Do you regularly lift or carry a load as heavy as a 25 kg sack of flour or rice? & 1 yes; 2 no & \textbf{Adapted from} STEP module 6B: 'do you regularly have to lift or pull anything weighing at least ...' (item label verified in STEP Philippines file; the exact weight in STEP must be checked). \\
03 & Is your work mostly the same tasks, day after day? & 1 yes; 2 no & \textbf{Adapted from} STEP module 6B 'short, repetitive tasks' (frequency scale recast as yes/no); Autor-Handel 2013 routine-task item. \\
04 & Do you learn new things in your work each season? & 1 yes; 2 no & \textbf{Adapted from} STEP module 6B 'learning new things' (frequency scale recast as yes/no). \\
05 & Who decides how you do your work? & 1 I decide; 2 shared (for example with a union or partner); 3 someone else (owner, union, customer) & \textbf{Adapted from} Apablaza et al. 2026, Table 4: 'autonomy to decide how and when tasks are performed'; STEP autonomy index. \\
06 & Do you deal directly with customers or pilgrims on every working day? & 1 yes; 2 no & \textbf{Adapted from} Autor-Handel 2013 PDII items on face-to-face contact with customers or clients (share of day recast as yes/no). \\
07 & How many people do you regularly direct or supervise? & whole number (0 if none) & \textbf{Adapted from} Autor-Handel 2013 'proportion of workday managing or supervising others'; STEP 'supervise others' (recast as a count). \\
08 & How long did it take YOU to learn your work well? & number of months or seasons (enumerator codes to bins later) & \textbf{Adapted from} STEP module 6B 'how long would it take someone to learn to do this work well' (recast as the respondent's own experience). \\
\bottomrule
\end{longtable}
}

\section*{6. Questions about general skills and papers held}
These sixteen questions cover reading, counting, languages, phones, courses, and licences.
They matter because some jobs need a licence or a certificate, however many tasks a person already knows.
{\footnotesize
\begin{longtable}{p{0.7cm}p{4.5cm}p{2.8cm}p{6.6cm}}
\toprule
\textbf{No.} & \textbf{Question} & \textbf{Answers} & \textbf{Where it comes from} \\
\midrule
\endfirsthead
\toprule
\textbf{No.} & \textbf{Question} & \textbf{Answers} & \textbf{Where it comes from} \\
\midrule
\endhead
01 & In your work, do you read anything (notes, rate lists, tickets, messages, notices)? & 1 yes; 2 no & \textbf{Taken from} STEP module 6A: 'do you read anything at this work, including very short notes or instructions' (label verified). \\
02 & In your work, do you write anything (names, amounts, messages)? & 1 yes; 2 no & \textbf{Adapted from} Companion to the STEP module 6A reading item; STEP tests writing only through the language items. \\
03 & In your work, do you work out prices or costs? & 1 yes; 2 no & \textbf{Taken from} STEP module 6A: 'calculate prices or costs' (label verified). \\
04 & In your work, do you multiply, divide, or work out change or discounts? & 1 yes; 2 no & \textbf{Taken from} STEP module 6A: 'perform any other multiplication or division' and 'use or calculate fractions, decimals or percentages' (labels verified). \\
05 & Has not being able to read, write or do sums ever kept you from a job or a better job? & 1 yes; 2 no & \textbf{Taken from} STEP module 6A: 'has a lack of reading and writing skills ever kept you from ...' (label verified). \\
06 & For each language: can you speak it well enough to work with customers? (Hindi, Garhwali or local language, Nepali, English, other) & yes / no for each language & \textbf{Adapted from} Sengupta et al. 2021 language skills (speaking, listening); the scale is dropped. \\
07 & For each language: can you read and write it? & yes / no for each language & \textbf{Taken from} STEP module 7 asks read/write for each language (labels verified). \\
08 & What phone do you use? (none, basic phone, smartphone; own or shared) & list & \textbf{Adapted from} Sengupta et al. 2021 'mobile phone usage', 'internet usage'; Pal et al. 2026 'use internet on phone'; project pilot variable digital\_smartphone\_use. \\
09 & What do you use the phone for? (calls, messages or WhatsApp, internet, search or maps, digital payments) & multiple answers & \textbf{Adapted from} Pal et al. 2026 generic-skills block; project pilot variable digital\_payment\_use. \\
10 & In the last 12 months, have you used a computer or tablet? If yes, for what? & yes / no; then list & \textbf{Taken from} STEP module 6B computer-use items (labels verified); Pal et al. 2026 (about 33\% of workers used computers). \\
11 & Have you ever completed a training course or apprenticeship (not school)? & 1 yes; 2 no & \textbf{Taken from} STEP module 2: 'did you participate in any training or program, not as part of a formal education' (label verified). \\
12 & What kind of course was it? (driving, cooking or hospitality, guiding or trekking, electrical or technical, building trades, animal care, first aid or safety, computer, other) & multiple answers & \textbf{Adapted from} STEP module 2 course types (computer, basic skills, language, occupation-specific, personal development) recast to local trades. \\
13 & Did you receive a certificate? & 1 yes; 2 no & \textbf{Adapted from} NCO-2015 links some jobs to Qualification Packs with NSQF levels (for example electrician level 4, loader level 3). \\
14 & Who ran the course? (government scheme, NGO, private, union or association, employer) & single answer & \textbf{Adapted from} Pal et al. 2026 'access to skilling opportunities' section. \\
15 & Do you hold any of these? Please show it if you can. (two-wheeler licence, car licence, heavy vehicle licence, guide or trekking permit, food registration, ITI or skill certificate, first-aid certificate, security-guard certificate, none) & multiple answers; enumerator records seen / not seen & \textbf{Adapted from} NCO 8322.0501 (a valid driving licence is required); NCO qualification packs for electrician (level 4), gateman (level 4). \\
16 & How did you learn your main work? (family or inherited, watching others, on the job, on my own, a course) & multiple answers & \textbf{Adapted from} Pal et al. 2026: most workers learned by doing and from peers, only about one quarter had formal training first. \\
\bottomrule
\end{longtable}
}

\section*{7. The groups we will use later}
To see the big picture, we sort the 27 tasks into groups. We do this in two steps, general first.
\begin{itemize}[leftmargin=1.2em,itemsep=2pt]
\item \textbf{Three broad groups} (the same split used by Gathmann and Sch\"onberg, 2010): \emph{thinking and counting} (4 tasks), \emph{hands-on} (16 tasks), and \emph{dealing with people} (7 tasks).
\item \textbf{Five finer groups} (following Autor, Levy and Murnane, 2003, and Spitz-Oener, 2006). These also ask whether the work is repeated the same way each time or varies. The five are: hands-on varied (13 tasks), hands-on repeated (3), checking and counting repeated (3), judging varied (1), and dealing with people varied (7).
\end{itemize}
The three broad groups follow the published paper. The finer five-way sorting is our own judgement,
using the examples in the papers. We will record it before the survey goes to the field.

\section*{8. How solid is each question}
There are 52 questions in all. 28 are \textbf{taken from} a published survey or the national occupation list,
21 are \textbf{adapted} (the same idea, changed for spoken use or to remove a rating scale),
and 3 are \textbf{new}. The new ones are:
\begin{itemize}[leftmargin=1.2em,itemsep=2pt]
\item \textbf{t1\_02}: Carry a passenger (palki, dandi, kandi) \textit{Why it is new:} No NCO-2015 code for palki/dandi bearers (documented gap); nearest NCO 9331.0100 hand and pedal vehicle driver. Needs union check.
\item \textbf{t1\_13}: Judge whether the weather or route is safe and decide whether to go \textit{Why it is new:} No published counterpart. Idea anchored in the non-routine problem-solving task of Autor-Levy-Murnane 2003 and the 30-minute problem-solving item of Autor-Handel 2013.
\item \textbf{t2\_01}: In the busiest weeks, how many trips do you make in a day? (asked only if a carrying or animal task is coded 1) \textit{Why it is new:} Route-specific count. Replaces the STEP 1-10 physical-effort scale, which is a rating. Needs union check.
\end{itemize}

\section*{9. What is still to be checked before the survey}
\begin{itemize}[leftmargin=1.2em,itemsep=2pt]
\item \textbf{Talk to the unions} (pony, palki-kandi, jeep) to add or drop tasks. Palki and dandi work has no code in the national list.
\item \textbf{Check the exact weight and wording} of the lifting and repetition questions against the original STEP questionnaire. Our file only shows shortened labels.
\item \textbf{Ask workers in their own words first.} A short round (30 to 40 workers) asking ``what do you know how to do that helps you earn?'' follows the method of Sengupta and colleagues (2021). It will show if the task list is missing something.
\item \textbf{Make picture cards}, translate into Hindi, Garhwali and Nepali, and test with about 10 interviews.
\end{itemize}
\end{document}
